\documentclass[11pt]{article}

\usepackage[margin=1in]{geometry}
\usepackage{amsmath}
\usepackage{amssymb}
\usepackage{amsthm}

\newtheorem{theorem}{Theorem}
\newtheorem{lemma}[theorem]{Lemma}
\newtheorem{proposition}[theorem]{Proposition}
\newtheorem{corollary}[theorem]{Corollary}
\theoremstyle{definition}
\newtheorem{definition}[theorem]{Definition}
\newtheorem{conjecture}[theorem]{Conjecture}
\theoremstyle{remark}
\newtheorem{remark}[theorem]{Remark}

\title{Disproof of conjecture 00000001070:\\
the three-fold sum threshold in $\mathbb{F}_p$}
\author{Submission for conjecture 00000001070}
\date{2026}

\begin{document}
\maketitle

\begin{abstract}
Conjecture 00000001070 asserts that the smallest cardinality $|A|$ of a set
$A \subseteq \mathbb{F}_p$ with three-fold sum $3A = \mathbb{F}_p$ equals
$\lceil (p+1)/3 \rceil + 1$. We show that this is \textbf{false}. The
three-element set $A = \{0,1,3\} \subseteq \mathbb{F}_7$ already satisfies
$3A = \mathbb{F}_7$, while the conjectured expression evaluates to
$\lceil 8/3 \rceil + 1 = 4 \neq 3$. We also explain the arithmetic slip: the
value $\lceil (p+2)/3 \rceil$, obtained from Cauchy--Davenport, is the
\textbf{universal} threshold (every $k$-element set covers $\mathbb{F}_p$ when
$k \ge \lceil (p+2)/3 \rceil$), not the \textbf{existential} minimum stated in
the conjecture; the latter is far smaller and is known only through additive
basis theory. The conjecture is wrong under either reading.
\end{abstract}

\section{The conjecture}

\begin{definition}[Three-fold sumset]
Let $p$ be prime and let $A \subseteq \mathbb{F}_p$ be nonempty. The
\emph{three-fold sumset} is
\[
  3A \;=\; A + A + A \;=\; \{\, a + b + c : a, b, c \in A \,\},
\]
where repetitions of summands are allowed, and addition is in
$\mathbb{F}_p = \mathbb{Z}/p\mathbb{Z}$.
\end{definition}

\begin{conjecture}[as filed]\label{conj}
The smallest $|A|$ with $3A = \mathbb{F}_p$ is
\[
  \Bigl\lceil \frac{p+1}{3} \Bigr\rceil + 1 .
\]
\end{conjecture}

\section{The counterexample at $p=7$}

\begin{theorem}\label{thm:main}
Conjecture~\ref{conj} is false. Explicitly,
\[
  A = \{0,1,3\} \subseteq \mathbb{F}_7,
  \qquad |A| = 3,
  \qquad 3A = \mathbb{F}_7,
\]
while
\[
  \Bigl\lceil \frac{7+1}{3} \Bigr\rceil + 1 = \lceil 8/3 \rceil + 1 = 3 + 1 = 4 \neq 3 .
\]
Thus the conjectured expression overestimates the smallest $|A|$ already at
$p = 7$.
\end{theorem}

\begin{proof}
The set $A$ has three distinct elements. Enumerating the $3 \times 3 \times 3 =
27$ ordered triples, the distinct three-fold sums are
\[
\begin{array}{c|c}
y & \text{representation as } a+b+c \\ \hline
0 & 0+0+0 \\
1 & 0+0+1 \\
2 & 0+1+1 \\
3 & 0+0+3 \\
4 & 0+1+3 \\
5 & 1+1+3 \\
6 & 0+3+3
\end{array}
\]
so $3A = \{0,1,2,3,4,5,6\} = \mathbb{F}_7$. The arithmetic value of the
conjectured formula is $4$, and $3 < 4$.
\end{proof}

\begin{remark}
No two-element subset of $\mathbb{F}_7$ has three-fold sum $\mathbb{F}_7$;
the $21 = \binom{7}{2}$ two-element subsets each produce at most $4$ distinct
sums. Hence the true minimum at $p=7$ is exactly $3$, and the conjectured
value $4$ is not merely unattained but strictly larger than the truth.
\end{remark}

\section{Why the formula is wrong: two different thresholds}

The conjectured expression is motivated by Cauchy--Davenport, but it answers a
different question from the one that is stated. We separate the two.

\subsection{The universal threshold}

\begin{theorem}[Cauchy--Davenport]\label{thm:cd}
For nonempty $A, B \subseteq \mathbb{F}_p$ with $p$ prime,
\[
  |A + B| \;\ge\; \min\bigl(p,\; |A| + |B| - 1\bigr).
\]
\end{theorem}

Iterating Theorem~\ref{thm:cd} twice gives $|3A| \ge \min(p,\, 3|A| - 2)$:
\begin{align*}
  |2A| &\ge \min(p,\, 2|A| - 1), \\
  |3A| = |A + 2A| &\ge \min\bigl(p,\, |A| + \min(p, 2|A|-1) - 1\bigr)
                    \;\ge\; \min(p,\, 3|A| - 2).
\end{align*}
Consequently, if $3|A| - 2 \ge p$ then $|3A| \ge p$, so $3A = \mathbb{F}_p$
for \emph{every} set $A$ of that size. Writing $k = |A|$, the least $k$ with
$3k - 2 \ge p$ is
\[
  k_{\mathrm{univ}}(p) \;=\; \Bigl\lceil \frac{p+2}{3} \Bigr\rceil .
\]
This is sharp for the universal statement: for $k < \lceil (p+2)/3 \rceil$ the
interval $A = \{0,1,\dots,k-1\}$ has $3A = \{0,1,\dots,3k-3\}$ of size
$3k - 2 < p$, hence does not cover $\mathbb{F}_p$.

\begin{proposition}\label{prop:conj-vs-univ}
The conjectured value $\lceil (p+1)/3 \rceil + 1$ equals the universal
threshold $\lceil (p+2)/3 \rceil$ if and only if $p \equiv 2 \pmod 3$. It is
too large by exactly $1$ for $p = 3$ and for every prime $p \equiv 1 \pmod 3$;
for $p \le 30$ these are
\[
  p = 3,\ 7,\ 13,\ 19 .
\]
\end{proposition}

\begin{proof}
Write the remainder of $p$ modulo $3$. If $p = 3m + 2$, then
$\lceil (p+1)/3 \rceil + 1 = m + 2 = \lceil (p+2)/3 \rceil$. If $p = 3m+1$,
then $\lceil (p+1)/3 \rceil + 1 = m + 2$ while $\lceil (p+2)/3 \rceil = m+1$,
a difference of $1$. If $p = 3$ then $\lceil 4/3 \rceil + 1 = 3$ while
$\lceil 5/3 \rceil = 2$, again a difference of $1$.
\end{proof}

So even if the intended quantity were the universal threshold, the conjecture
would be wrong at $p = 3$ and at every prime $p \equiv 1 \pmod 3$.

\subsection{The existential minimum stated in the conjecture}

Conjecture~\ref{conj} literally asks for the \emph{smallest} $|A|$ for which
\emph{some} $A$ covers $\mathbb{F}_p$. This existential minimum is much
smaller than the universal threshold, because a $k$-element set can produce
\[
  \Bigl|\{\,a+b+c : a,b,c \in A\,\}\Bigr| \;\le\; \binom{k+2}{3}
  \;=\; \frac{k(k+1)(k+2)}{6}
\]
distinct sums (choosing an unordered multiset of three summands). Covering all
$p$ residues therefore forces $\binom{k+2}{3} \ge p$, so
$k \gtrsim (6p)^{1/3}$; the true threshold grows like $p^{1/3}$, not linearly
in $p$. Exhaustive search for all primes $p \le 30$ gives the following
values, where ``true'' is the smallest $|A|$ with some $3A = \mathbb{F}_p$.

\begin{center}
\begin{tabular}{r|r|r|r}
$p$ & true (existential) & $\lceil (p+1)/3 \rceil + 1$ (conjecture)
    & $\lceil (p+2)/3 \rceil$ (universal) \\ \hline
 2 & 2 &  2 &  2 \\
 3 & 2 &  3 &  2 \\
 5 & 3 &  3 &  3 \\
 7 & 3 &  4 &  3 \\
11 & 4 &  5 &  5 \\
13 & 4 &  6 &  5 \\
17 & 5 &  7 &  7 \\
19 & 5 &  8 &  7 \\
23 & 5 &  9 &  9 \\
29 & 6 & 11 & 11
\end{tabular}
\end{center}

\begin{remark}
For $p \le 30$ the conjectured value matches the true existential minimum only
at $p = 2$ and $p = 5$; at $p = 3, 7, 11, 13, 17, 19, 23, 29$ it is strictly
too large. The universal threshold likewise fails to be the existential
minimum from $p = 11$ onward (for instance $A = \{0,1,2,4\}$ of size $4$
covers $\mathbb{F}_{11}$, while $\lceil 13/3 \rceil = 5$). Thus the
conjectured formula is not the answer to either question.
\end{remark}

\section{Formal verification}

The counterexample of Theorem~\ref{thm:main} is machine-checked in the
accompanying Lean~4 project (\texttt{lean4/}). The field $\mathbb{F}_7$ is
modelled as \texttt{Fin 7} with its cyclic addition; the three-fold sum is a
decidable predicate
\[
  \texttt{IsThreefoldSum}\ A\ y \;:\Longleftrightarrow\;
  \exists a\,b\,c : \texttt{Fin }7,\ a \in A \wedge b \in A \wedge c \in A
  \wedge a + b + c = y ,
\]
and the project proves
\[
  \forall y : \texttt{Fin }7,\ \texttt{IsThreefoldSum}\ \{0,1,3\}\ y,
  \qquad |\{0,1,3\}| = 3,
  \qquad \lceil 8/3 \rceil + 1 = 4 \neq 3 .
\]
Only core Lean (\texttt{import Std}, no Mathlib, no \texttt{ZMod}, no
\texttt{sorry}) is used.

\section{Conclusion}

Conjecture 00000001070 is \textbf{false}. The set $A = \{0,1,3\}$ shows
$3A = \mathbb{F}_7$ with $|A| = 3$, whereas the conjectured formula returns
$4$. The underlying identity $\lceil (p+2)/3 \rceil = \lceil (p+1)/3 \rceil +
1$ has the wrong parity of rounding: the two sides agree only when
$p \equiv 2 \pmod 3$, and differ by one at $p = 3$ and all primes
$p \equiv 1 \pmod 3$. Moreover, neither expression is the existential minimum
that the conjecture states; that minimum grows like $p^{1/3}$.

\end{document}
